% Image credits for the photographs and AI illustrations of Book 1 (Primary and Middle School).
% Machine-readable license records live in images/book1/CREDITS.md; keep the
% two files in sync. The Book 1 writer owns this file: add one \item per
% photograph (inside the itemize, by chapter) as photographs land.
% Arabic edition: keep this file in sync with frontmatter/image-credits.tex
% (one \item per item there). Personal names are transliterated, with the
% original spelling in parentheses (the credit itself); institutions,
% datasets, works and licences stay verbatim; only the gloss is translated.
\chapter*{مصادر الصور}

الرسوم والخرائط والمخططات في هذا الكتاب أصلية؛ وما رُسم منها من بيانات
منشورة يذكر مصدره في التعليق المرافق أو في هذه القائمة. أما الصور
الفوتوغرافية فمستعملة بمقتضى رخصها الحرة، مع الشكر لأصحابها؛ وروابط المصادر
الكاملة وروابط الرخص لكل صورة مدرجة في \texttt{images/book1/CREDITS.md} في
مستودع الكتاب.

\bigskip
وإلى جانب الصور الفوتوغرافية والرسوم الأصلية، بعض الأشكال رسوم توضيحية
مولَّدة بالذكاء الاصطناعي، أُنتجت بأداة توليد الصور من OpenAI انطلاقًا من
أوامر كتبها محررو الكتاب، ورُوجعت دقتها العلمية قبل إدراجها. وهي لا تُظهر
إلا مشاهد عامة، ولا تمثّل أبدًا مكانًا حقيقيًا أو شخصًا حقيقيًا أو عيّنة
حقيقية. والأمر الكامل لكل صورة مولَّدة مسجَّل في
\texttt{images/book1/ai/PROMPTS.md} في المستودع.

\bigskip
\noindent الصور الفوتوغرافية والبيانات، فصلًا فصلًا:
\begin{itemize}\small
  \item الفصل~\ref{ch:g2:water-on-earth} (الماء على الأرض): الأرض كما رُئيت من \babelsublr{Apollo~17}، في 7~ديسمبر 1972 (\babelsublr{\emph{The Blue Marble}})؛ \babelsublr{NASA}، ملك عام (عبر \babelsublr{Wikimedia Commons}). والمربعات المئة محسوبة من تضاريس \babelsublr{NOAA ETOPO1}؛ والدلاء المئة من جدول \babelsublr{U.S. Geological Survey Water Science School} (عن شيكلومانوف (\babelsublr{Shiklomanov})، 1993).
  \item الفصل~\ref{ch:g3:stones-pebbles-sand} (الحجارة والحصى والرمل): جدار من حجارة الغرانيت المرصوفة بلا ملاط؛ صورة التقطها توني أتكين (\babelsublr{Tony Atkin})، \babelsublr{Geograph}، \babelsublr{CC BY-SA 2.0} (عبر \babelsublr{Wikimedia Commons}). وسقوف من ألواح الحجر في سويسرا؛ صورة التقطها مايكل ك.~رايغل (\babelsublr{Michael C. Rygel})، \babelsublr{CC BY-SA 3.0} (عبر \babelsublr{Wikimedia Commons}).
  \item الفصل~\ref{ch:g3:continents-and-oceans} (مجسم الكرة الأرضية): المجسمان وخريطة العالم مرسومة من سواحل \babelsublr{Natural Earth} (ملك عام).
  \item الفصل~\ref{ch:g4:clouds-rain-snow} (السحب والمطر والثلج): سحب ركامية فوق كانساس (س.~كورفيدي، \babelsublr{S.~Corfidi})، وبقع من السحب الريشية (ر.~ف.~كريسغ، \babelsublr{R.~F.~Kresge})، وطبقة من السحب الطبقية حول قمم جبلية، وعاصفة رعدية فوق صحراء أريزونا (مجموعة ب.~هوك، \babelsublr{B.~Hooke})، وحبة برد كبيرة جدًا؛ كلها من \babelsublr{NOAA Photo Library}، ملك عام (عبر \babelsublr{Wikimedia Commons}).
  \item الفصل~\ref{ch:g4:volcanoes-earthquakes-first-look} (البراكين والزلازل): حمم تدخل المحيط عند كيلاويا، \babelsublr{USGS Hawaiian Volcano Observatory}؛ وجبل سانت هيلينز في 18~مايو 1980، صورة التقطها أوستن بوست (\babelsublr{Austin Post})، \babelsublr{USGS}؛ وبيت انقلع عن أساسه في زلزال لوما بريتا عام 1989، صورة التقطها ج.~ك.~ناكاتا (\babelsublr{J.~K.~Nakata})، \babelsublr{USGS}؛ كلها ملك عام (عبر \babelsublr{Wikimedia Commons}). وتستعمل الخريطة قائمة البراكين الهولوسينية في \babelsublr{Smithsonian Global Volcanism Program}، وفهرس الزلازل من \babelsublr{U.S. Geological Survey}.
  \item الفصل~\ref{ch:g4:world-physical-map} (الخريطة الطبيعية): الهيمالايا والأمازون كما صُوّرا من محطة الفضاء الدولية، \babelsublr{NASA}، ملك عام (عبر \babelsublr{Wikimedia Commons}). والخريطة الطبيعية مرسومة من تضاريس \babelsublr{NOAA ETOPO1}، مع الأنهار والبحيرات والسواحل من \babelsublr{Natural Earth} (ملك عام).
  \item الفصل~\ref{ch:g5:water-shapes-the-land} (الماء ينحت الأرض): نهر ماموري كما صُوّر من محطة الفضاء الدولية (\babelsublr{ISS007-E-10797})، ودلتا النيل كما رآها جهاز \babelsublr{MODIS} على القمر الصناعي \babelsublr{Terra}، \babelsublr{NASA}، ملك عام (عبر \babelsublr{Wikimedia Commons}).
  \item الفصل~\ref{ch:g5:weather-and-climate} (الطقس والمناخ): صحراء سونورا، \babelsublr{U.S. National Park Service}، \babelsublr{Saguaro National Park}، ملك عام (عبر \babelsublr{Wikimedia Commons}). والمخططات المناخية مرسومة من \babelsublr{NOAA NCEI U.S. Climate Normals 1991--2020} (ملك عام).
  \item الفصل~\ref{ch:g5:natural-hazards} (الأخطار الطبيعية): عين الإعصار إيزابيل كما صوّرها من محطة الفضاء الدولية إد لو (\babelsublr{Ed Lu})، \babelsublr{NASA}، ملك عام؛ وانزلاق التربة عام 2014 في ولاية واشنطن، مارك ريد (\babelsublr{Mark Reid})، \babelsublr{U.S. Geological Survey}، ملك عام (كلتاهما عبر \babelsublr{Wikimedia Commons}).
  \item الفصل~\ref{ch:g6:earth-a-planet} (الأرض، كوكب): الزهرة من \babelsublr{Mariner~10} (\babelsublr{NASA/JPL-Caltech})؛ والأرض من \babelsublr{Apollo~17} (\babelsublr{NASA})؛ و\babelsublr{\emph{Earthrise}} من \babelsublr{Apollo~8} (\babelsublr{NASA}، بيل أندرز، \babelsublr{Bill Anders})؛ والمريخ والمشتري من تلسكوب هابل الفضائي (\babelsublr{NASA}، \babelsublr{ESA}، \babelsublr{the Hubble Heritage Team} و \babelsublr{A.~Simon})؛ كلها ملك عام (عبر \babelsublr{Wikimedia Commons}). و\babelsublr{\emph{Pale Blue Dot}} كما رآها \babelsublr{Voyager~1} عام 1990، نسخة 2020 (\babelsublr{NASA/JPL-Caltech}، \babelsublr{PIA23645}، ملك عام). وبيانات الكواكب من \babelsublr{NASA planetary fact sheet}.
  \item الفصل~\ref{ch:g6:latitude-longitude} (العرض والطول): ساعة هاريسون \babelsublr{H4} في \babelsublr{National Maritime Museum}، غرينتش، صورة التقطها \babelsublr{Colonel Warden}، \babelsublr{CC BY-SA 3.0} (عبر \babelsublr{Wikimedia Commons}). والسواحل والمدن من \babelsublr{Natural Earth} (ملك عام).
  \item الفصل~\ref{ch:g6:rocks-and-minerals} (الصخور والمعادن): الغرانيت، والبازلت، والحجر الجيري الأحفوري، والحجر الرملي، والرخام، والسجّيل، وبلورة كوارتز، ورقاقة من ميكا المسكوفيت، صور التقطها جيمس سانت جون (\babelsublr{James St.~John})، \babelsublr{CC BY 2.0}؛ وسقف من الأردواز، صورة التقطها أويفيند هولمستاد (\babelsublr{Øyvind Holmstad})، \babelsublr{CC0} (كلها عبر \babelsublr{Wikimedia Commons}).
  \item الفصل~\ref{ch:g6:inside-the-earth} (باطن الأرض): إنغه ليمان (\babelsublr{Inge Lehmann}) عام 1932، صورة التقطها إيفن نويهاوس (\babelsublr{Even Neuhaus})، \babelsublr{Royal Library of Denmark}، ملك عام (عبر \babelsublr{Wikimedia Commons}). والطبقات والكثافات وسرعات الموجات مرسومة من \babelsublr{Preliminary Reference Earth Model} (أ.~م.~دزيفونسكي ود.~ل.~أندرسون، \babelsublr{A.~M. Dziewonski and D.~L. Anderson}، 1981).
  \item الفصل~\ref{ch:g6:water-underground} (الماء تحت الأرض): \babelsublr{Big Spring} في ميزوري، صورة التقطها \babelsublr{Johnston9494}، \babelsublr{CC BY-SA 4.0} (عبر \babelsublr{Wikimedia Commons}). والمساميات عن ر.~ك.~هيث (\babelsublr{R.~C. Heath})، \babelsublr{\emph{Basic Ground-Water Hydrology}}، \babelsublr{U.S. Geological Survey} (1983)؛ وحصص الماء العذب من \babelsublr{USGS Water Science School}.
  \item الفصل~\ref{ch:g7:plate-tectonics} (تكتونية الصفائح): القارات الجنوبية وقد أُعيد جمعها، مع نطاقات الأحافير الأربع: \babelsublr{\emph{Cynognathus}}، \babelsublr{\emph{Lystrosaurus}}، \babelsublr{\emph{Mesosaurus}}، \babelsublr{\emph{Glossopteris}}؛ \babelsublr{U.S. Geological Survey}، \babelsublr{\emph{This Dynamic Earth}}، ملك عام.
  \item الفصل~\ref{ch:g7:plate-tectonics}: صدع سان أندرياس عبر سهل كاريزو؛ صورة التقطها روبرت إ.~والاس (\babelsublr{Robert E. Wallace})، \babelsublr{U.S. Geological Survey}، ملك عام.
  \item الفصل~\ref{ch:g7:plate-tectonics}: حدود الصفائح عن ب.~بيرد (\babelsublr{P.~Bird}) (2003)، بصيغة \babelsublr{GeoJSON} من إعداد ه.~أهلينيوس (\babelsublr{H.~Ahlenius}) (\babelsublr{ODC-By 1.0})؛ والسواحل من \babelsublr{Natural Earth} (ملك عام)؛ والزلازل من فهرس \babelsublr{U.S. Geological Survey}؛ والبراكين من \babelsublr{Smithsonian Global Volcanism Program}.
  \item الفصل~\ref{ch:g7:volcanoes} (البراكين): جبل سانت هيلينز من \babelsublr{Johnston Ridge} في 17~مايو و10~سبتمبر 1980، صورتان التقطهما هاري غليكن (\babelsublr{Harry Glicken})، \babelsublr{U.S. Geological Survey}؛ ونافورة حمم عند كيلاويا عام 1983، وماونا لوا من الجو عام 1985، صورتان التقطهما ج.~د.~غريغز (\babelsublr{J.~D. Griggs})، \babelsublr{U.S. Geological Survey}؛ وعمود ثوران جبل بيناتوبو في 12~يونيو 1991، صورة التقطها ديف هارلو (\babelsublr{Dave Harlow})، \babelsublr{U.S. Geological Survey}؛ كلها ملك عام (عبر \babelsublr{Wikimedia Commons}). وتستعمل الخريطة والأعداد قائمة البراكين الهولوسينية في \babelsublr{Smithsonian Global Volcanism Program}؛ وحدود الصفائح عن ب.~بيرد (\babelsublr{P.~Bird}) (2003).
  \item الفصل~\ref{ch:g7:earthquakes} (الزلازل): قرية قرب ساحل سومطرة بعد تسونامي 26~ديسمبر 2004، صورة من \babelsublr{U.S. Navy} التقطها فيليب أ.~ماكدانيال (\babelsublr{Philip A. McDaniel})؛ وحافات صدوع زلزال بورا بيك عام 1983، صورة التقطها ر.~إ.~والاس (\babelsublr{R.~E. Wallace})، \babelsublr{U.S. Geological Survey}؛ والطريق المنهار في بحيرة هيبغن عام 1959، صورة التقطها إ.~ج.~ويتكايند (\babelsublr{I.~J. Witkind})، \babelsublr{U.S. Geological Survey}؛ كلها ملك عام (عبر \babelsublr{Wikimedia Commons}). وتستعمل الأعداد والأعماق فهرس الزلازل من \babelsublr{U.S. Geological Survey}؛ وجدول الشدة مختصر من سلّم ميركالي المعدَّل كما تنشره \babelsublr{U.S. Geological Survey}.
  \item الفصل~\ref{ch:g7:sediments-and-fossils} (الرواسب والأحافير): الأخدود العظيم من \babelsublr{Hermits Rest}، صورة من \babelsublr{NPS} التقطتها كريستن م.~كالدون (\babelsublr{Kristen M. Caldon})، ملك عام؛ وأمونيت وثلاثي فصوص وعلامات تموّج في \babelsublr{Capitol Reef National Park}، صور التقطها جيمس سانت جون (\babelsublr{James St.~John})، \babelsublr{CC BY 2.0}؛ وآثار أقدام ديناصورات في \babelsublr{Dinosaur Ridge}، صورة التقطها دكستر بيركنز (\babelsublr{Dexter Perkins})، \babelsublr{CC0} (كلها عبر \babelsublr{Wikimedia Commons}).
  \item الفصل~\ref{ch:g7:geological-time} (أعمار الأرض): جيمس هاتون (\babelsublr{James Hutton})، صورة رسمها هنري ريبيرن (\babelsublr{Henry Raeburn})، ملك عام (عبر \babelsublr{Wikimedia Commons}). وأعمار السلّم الزمني وألوانه هي أعمار \babelsublr{International Chronostratigraphic Chart} وألوانه (\babelsublr{International Commission on Stratigraphy}).
  \item الفصل~\ref{ch:g8:the-atmosphere} (الغلاف الجوي): إطلاق بالون أرصاد جوية في توبيكا، كانساس، صورة التقطها ستيفن كورفيدي (\babelsublr{Stephen Corfidi})، \babelsublr{NOAA/National Weather Service}؛ وحافة الأرض عند الغروب من محطة الفضاء الدولية، \babelsublr{NASA}، \babelsublr{ISS Expedition~23}؛ كلتاهما ملك عام (عبر \babelsublr{Wikimedia Commons}). والمقاطع الرأسية محسوبة من \babelsublr{\emph{U.S. Standard Atmosphere, 1976}} (\babelsublr{NOAA} و \babelsublr{NASA} و \babelsublr{U.S. Air Force}).
  \item الفصل~\ref{ch:g8:winds-global-circulation} (الرياح والدورة العامة): الأمريكتان كما رآهما القمر الصناعي \babelsublr{GOES-16} في 26~يوليو 2022، \babelsublr{NOAA}، ملك عام (عبر \babelsublr{Wikimedia Commons}). وسواحل خريطة الرياح التجارية من \babelsublr{Natural Earth} (ملك عام).
  \item الفصل~\ref{ch:g8:fronts-and-weather-maps} (الجبهات وخرائط الطقس): منخفض جوي قبالة آيسلندا كما رآه القمر الصناعي \babelsublr{Terra}، \babelsublr{NASA/GSFC}، \babelsublr{MODIS Rapid Response Team}، جاك ديكلوتر (\babelsublr{Jacques Descloitres})، ملك عام (عبر \babelsublr{Wikimedia Commons}). وسواحل خريطة الكتل الهوائية من \babelsublr{Natural Earth} (ملك عام).
  \item الفصل~\ref{ch:g8:the-oceans} (المحيطات): السفينة \babelsublr{HMS \emph{Challenger}} تستعد لسبر العمق عام 1872، نقش من تقارير بعثة \babelsublr{Challenger} أعيد نشره في و.~أ.~هيردمان (\babelsublr{W.~A. Herdman})، \babelsublr{\emph{Founders of Oceanography}} (1923)؛ وتيار الخليج كما رآه جهاز \babelsublr{MODIS} على القمر الصناعي \babelsublr{Aqua}، \babelsublr{NASA Earth Observatory}؛ كلتاهما ملك عام (عبر \babelsublr{Wikimedia Commons}). ومقاطع قاع المحيط مرسومة من تضاريس \babelsublr{NOAA ETOPO1}؛ وسواحل خريطتي الدوامات والعوامات من \babelsublr{Natural Earth} (ملك عام).
  \item الفصل~\ref{ch:g8:ice-on-earth} (الجليد على الأرض): نهر بيدرسن الجليدي عامي 1990 و2016، صورتان التقطهما م.~تيترو (\babelsublr{M.~Tetreau}) ود.~كورتز (\babelsublr{D.~Kurtz})، \babelsublr{U.S. National Park Service}؛ وحافة جرف جليدي قبالة شبه الجزيرة القطبية الجنوبية، \babelsublr{NASA}، \babelsublr{Operation IceBridge} (م.-ج.~فينياس، \babelsublr{M.-J. Viñas})؛ كلها ملك عام (عبر \babelsublr{Wikimedia Commons}). ويستعمل مخطط الجليد البحري \babelsublr{NSIDC Sea Ice Index}، الإصدار~4 (ف.~فيترر (\babelsublr{F.~Fetterer}) وآخرون، 2025، \babelsublr{doi:10.7265/a98x-0f50}).
  \item الفصل~\ref{ch:g9:climate-past-and-present} (مناخ الأرض بين الماضي والحاضر): مقطع من لب الجليد \babelsublr{GISP2} من غرينلاند، \babelsublr{NOAA Paleoclimatology}، ملك عام (عبر \babelsublr{Wikimedia Commons}). والسجلات مرسومة من \babelsublr{NASA GISTEMP} الإصدار~4، وسجل ثاني أكسيد الكربون في ماونا لوا (\babelsublr{NOAA Global Monitoring Laboratory} و \babelsublr{Scripps Institution of Oceanography})، وحرارة \babelsublr{EPICA Dome~C} (ج.~جوزيل (\babelsublr{J.~Jouzel}) وآخرون، 2007)، والسجل المركّب لثاني أكسيد الكربون في ألباب الجليد القطبية الجنوبية (ب.~بيريتر (\babelsublr{B.~Bereiter}) وآخرون، 2015)، ومستوى سطح البحر من \babelsublr{NOAA Laboratory for Satellite Altimetry}؛ وعبارات التقييم هي عبارات \babelsublr{IPCC Sixth Assessment Report}، \babelsublr{Working Group~I}، \babelsublr{Summary for Policymakers} (2021).
  \item الفصل~\ref{ch:g9:natural-risks} (العيش مع المخاطر الطبيعية): خريطة أخطار نيفادو دل رويث، عن ت.~ل.~رايت (\babelsublr{T.~L. Wright}) وت.~ك.~بيرسون (\babelsublr{T.~C. Pierson})، \babelsublr{\emph{Living with Volcanoes}}، \babelsublr{U.S. Geological Survey Circular 1073} (1992)؛ وعوامة إنذار من التسونامي يفحصها طاقم \babelsublr{USCGC \emph{William Hart}}، \babelsublr{U.S. Coast Guard}، 2021؛ كلتاهما ملك عام (عبر \babelsublr{Wikimedia Commons}).
  \item الفصل~\ref{ch:g9:earth-resources} (موارد من الأرض): منجم تشوكيكاماتا كما رآه \babelsublr{ASTER}، \babelsublr{NASA/METI/AIST/Japan Space Systems} و \babelsublr{the U.S./Japan ASTER Science Team}، 2007، ملك عام؛ وطبقة فحم في مقطع طريق في مقاطعة بلمونت، أوهايو، صورة التقطها جيمس سانت جون (\babelsublr{James St.~John})، \babelsublr{CC BY 2.0} (كلتاهما عبر \babelsublr{Wikimedia Commons}). ومعاملات التركيز من د.~بيركنز (\babelsublr{D.~Perkins}) وآخرين، \babelsublr{\emph{Mineralogy}}، ومن س.~إيرل (\babelsublr{S.~Earle})، \babelsublr{\emph{Environmental Geology}}.
  \item الفصل~\ref{ch:g9:landforms} (الأشكال التضاريسية): الأخوات السبع (\babelsublr{Seven Sisters})، صورة التقطها أريلد فوغن (\babelsublr{Arild Vågen})، 2017، \babelsublr{CC BY-SA 4.0}؛ وكثبان هلالية على ساحل جنوب البرازيل، صورة التقطها رائد فضاء من \babelsublr{NASA} من محطة الفضاء الدولية، 2017، ملك عام؛ وهوابط كلسية في \babelsublr{Carlsbad Cavern}، صورة من \babelsublr{National Park Service} التقطها دانيال لايفهايت (\babelsublr{Daniel Leifheit})، 2018، ملك عام (كلها عبر \babelsublr{Wikimedia Commons}).
  \item الفصل~\ref{ch:g9:maps-of-the-earth} (خرائط الأرض كلها): دلتا نهر لينا، صورة \babelsublr{Landsat} من عام 2000، \babelsublr{USGS EROS Data Center} عبر \babelsublr{NASA Visible Earth}، ملك عام؛ وقمر صناعي \babelsublr{GPS}، تصوّر فني من \babelsublr{NASA}، ملك عام (كلتاهما عبر \babelsublr{Wikimedia Commons}). وسواحل خرائط العالم من \babelsublr{Natural Earth} (ملك عام).
\end{itemize}
